\section{Geometric invariants based on projections}


\paragraph*{Main notations and concepts.}
Let $\R^d$ be the Euclidean $d$-dimensional space and $\Z^d$ its
subset of points with integer coordinates. For any integer $m$, $0 \Le
m \Le d$, we denote by $\mathcal{I}_m$ the set of subsets of
$\{1,\dots,d\}$ with cardinal $m$. For any $J \in \mathcal{I}_m$, the
map $\pi_J: \R^d \to \R^m$ is the orthogonal projection onto the
subspace with coordinates in $J$. Let $p$ be any point of $\R^d$ and
let $r$ be a positive value which defines the radius of the local
analysis. We write $B(p,r)$ the $d$-dimensional ball of center $p$ and
radius $r$. The unit $m$-dimensional ball centered on $0$ is denoted
by $B^m$, with volume $\omega_m$. If $T$ is a subset of $\R^d$, we write
$B_T(p,r)$ for $B(p,r) \cap T$. The following measures associate a
local value to a set through projections onto sub-axes.

\begin{definition}[local measures]
  \label{def-local-measures}
  Let $\HH^m$ be the $m$-dimensional
  Hausdorff measure. For any subset $Y \subset \R^d$, $p \in Y$, $r>0$, we define the local measures
  \begin{align}
    V_J(Y,p,r) &:= \HH^m\!\left( \pi_J\bigl(B_Y(p,r)\bigr) \right),
    \quad
    && & W_m(Y,p,r) & := \sum_{J \in \mathcal{I}_m} V_J(Y,p,r)^2.
  \end{align}
\end{definition}

We will establish several stability results for these measures:
\begin{enumerate}
\item For $T$ an arbirary $k$-dimensional plane, $W_m(T,p,r)$ is $0$ if
  $m>k$, otherwise is a constant times the square of the volume of the
  $m$-ball, independently of the position and the orientation of $T$ in $\R^d$.
\item For $X$ a smooth closed $k$-manifold, $W_m(X,p,r)$ is close to
  $W_m(T_p X,p,r)$ (for $T_p X$ the tangent plane to $X$ at $p$),
  hence also independent of the orientation of the tangent plane.
\item The same result hold for a covering discretization of $X$.
\item The quantities $V_J$ are also stable and related to the local
  geometry of $X$ or its discretization (tangent vector, normal).
\end{enumerate}
Clearly, these properties give mathematical tools to analyse the local
dimension of arbitrary lattice sets as well as some geometric
information (see section on experiments). The stability estimates established below show that these quantities are
robust under controlled geometric perturbations and discretization noise. We will also show the convergence of
these local measures for finer and finer digitizations.

We need a few more definitions. Let $X \subset \R^d$ be a $C^2$-smooth
submanifold of $\R^d$ of dimension $k$ without boundary. We further
suppose it has positive reach $\reach(X)\geq \rho>0$. For $h>0$ a
discretization step, let $Q_h:= \left[-\frac{h}{2},
  \frac{h}{2}\right]^d$. The \emph{$h$-covering discretization of $X$}
is $G_h(X) := \{ z \in h\Z^d~|~X \cap (z+Q_h) \neq \emptyset \}$, its
\emph{$h$-voxelization} is $U_h(X) := \cup_{z \in G_h(X)} (z + Q_h)$.

Our main result below states that the local measure of a digitized set
depends essentially on the dimension of its underlying submanifold,
with an error that tends to zero with the discretization
gridstep. This provides a tool for analyzing the local dimension of
digital data.
\begin{theorem}[Local approximation of $W_m$]
\label{thm:main}
There exist a constant $C=C(d,k)>0$ such that, for every $p \in X$ and $0<m\Le k$, if $0<r\Le \frac{1}{1+\sqrt{d}}\rho$ and $0<h\Le r$, then
\begin{align}\label{eq:main}
  \left| \frac{1}{r^{2m}}W_m(U_h(X),p,r) - \omega_m^2\binom{k}{m} \right|
  \Le  C \left( \frac{h}{r} + \frac{r}{\rho} \right).
\end{align}
\end{theorem}
\begin{proof}
  The proof relies on three theorems proved afterwards:
  \RefTheorem{thm:ball-invariants} shows that $W_m(B^k)$ is a constant
  depending on $k$ and $m$, and independent of the orientation or
  position of the unit $k$-ball $B^k$ in
  $\R^d$. \RefTheorem{thm:stabilite-projections} proves that both
  measures $V_j$ and $W_m$ are stable
  with respect to controlled variations of $X$. Finally
  \RefTheorem{thm:cubical-admissible-set} establishes that the
  $h$-voxelization of $X$ is an admissible controlled variation of
  $X$.

  More precisely, let $X_r:= X \cap B(p,r)$ and $Y_h:=U_h(X) \cap
  B(p,r)$. From \RefTheorem{thm:cubical-admissible-set}, we have $X_r
  \subset Y_h \subset (X_r)_{\varepsilon_h}$, with
  $\varepsilon_h=(1+2/\sqrt{3})\sqrt d\,h$. Applying \eqref{eq:stabilite-W} of
  \RefTheorem{thm:stabilite-projections} with $Y=Y_h$ and
  $\varepsilon=\varepsilon_h$ induces:
  \begin{equation*}
    \left| W_m(U_h(X),p,r)-\omega_m^2\binom{k}{m}r^{2m} \right|
    \Le C r^{2m-1} \left( \sqrt{d} h+\frac{r^2}{\rho} \right).
  \end{equation*}
  Dividing by $r^{2m}$ gives the result.
  \qed
\end{proof}
Balancing the two error terms suggests the scale
$r=\Theta(\sqrt{h\rho})$, for which the relative error is
$O(\sqrt{h/\rho})$. This result implies that the
rescaled measure $W_m$ characterizes the local dimension of the
underlying digitized manifold, independently of its orientation or its
embedding space.

\begin{figure}
  \begin{center}
    % for h = 0.5
\tikzset{math3d/.style={x= {(-0.39cm,-0.05cm)}, z={(0cm,0.4cm)},y={(0.2cm,-0.2cm)}}}
\tikzset{guide/.style={gray!65,densely dotted}}

\NewDocumentCommand{\DrawVoxel}{m m m m O{1.0}}{%
  \draw[fill=#4!10!white,draw=black,thin,opacity=#5] (#1-0.5,#2-0.5,#3+0.5) -- (#1+0.5,#2-0.5,#3+0.5) -- (#1+0.5,#2+0.5,#3+0.5) -- (#1-0.5,#2+0.5,#3+0.5) -- cycle;
  \draw[fill=#4!30!white,draw=black,thin,opacity=#5] (#1+0.5,#2-0.5,#3-0.5) -- (#1+0.5,#2-0.5,#3+0.5) -- (#1+0.5,#2+0.5,#3+0.5) -- (#1+0.5,#2+0.5,#3-0.5) -- cycle;
  \draw[fill=#4!40!white,draw=black,thin,opacity=#5] (#1-0.5,#2+0.5,#3-0.5) -- (#1+0.5,#2+0.5,#3-0.5) -- (#1+0.5,#2+0.5,#3+0.5) -- (#1-0.5,#2+0.5,#3+0.5) -- cycle;
}
\NewDocumentCommand{\DrawPoint}{m m m m O{1.0}}{%
  \draw[fill=#4,draw=black,thin,opacity=#5] (#1,#2,#3) node {\color{#4!50!white}$\bullet$};
}
  
%% \def\ProjX{-2.0}
%% \def\ProjY{-4.0}
%% \def\ProjZ{-2.0}
\def\ProjX{0.0}
\def\ProjY{0.0}
\def\ProjZ{0.0}
\NewDocumentCommand{\ProjectVoxel}{m m m m O{\ProjX} O{\ProjY} O{\ProjZ}}{%
  \draw[fill=blue!#4!white,draw=black,thin] (#1-0.5,#2-0.5,#7) -- (#1+0.5,#2-0.5,#7) -- (#1+0.5,#2+0.5,#7) -- (#1-0.5,#2+0.5,#7) -- cycle;
  \draw[fill=red!#4!white,draw=black,thin] (#5,#2-0.5,#3-0.5) -- (#5,#2-0.5,#3+0.5) -- (#5,#2+0.5,#3+0.5) -- (#5,#2+0.5,#3-0.5) -- cycle;
  \draw[fill=green!#4!white,draw=black,thin] (#1-0.5,#6,#3-0.5) -- (#1+0.5,#6,#3-0.5) -- (#1+0.5,#6,#3+0.5) -- (#1-0.5,#6,#3+0.5) -- cycle;
  \draw[fill=cyan!#4!white,draw=black,thin] (#1-0.5,#6,#7) -- (#1-0.5,#6+0.1,#7) -- (#1+0.5,#6+0.1,#7) -- (#1+0.5,#6,#7) -- (#1+0.5,#6,#7+0.1) -- (#1-0.5,#6,#7+0.1) --cycle;
  \draw[fill=magenta!#4!white,draw=black,thin] (#5,#2-0.5,#7) -- (#5+0.1,#2-0.5,#7) -- (#5+0.1,#2+0.5,#7) -- (#5,#2+0.5,#7) -- (#5,#2+0.5,#7+0.1) -- (#5,#2-0.5,#7+0.1) --cycle;
  \draw[fill=yellow!#4!white,draw=black,thin] (#5,#6,#3-0.5) -- (#5+0.1,#6,#3-0.5) -- (#5+0.1,#6,#3+0.5) -- (#5,#6,#3+0.5) -- (#5,#6+0.1,#3+0.5) -- (#5,#6+0.1,#3-0.5) --cycle;
  \draw[guide] (#1-0.5,#6,#7) -- (#1-0.5,#2-0.5,#7);
  \draw[guide] (#1+0.5,#6,#7) -- (#1+0.5,#2-0.5,#7);
  \draw[guide] (#1-0.5,#6,#7) -- (#1-0.5,#6,#3-0.5);
  \draw[guide] (#1+0.5,#6,#7) -- (#1+0.5,#6,#3-0.5);
  \draw[guide] (#5,#2-0.5,#7) -- (#1-0.5,#2-0.5,#7);
  \draw[guide] (#5,#2+0.5,#7) -- (#1-0.5,#2+0.5,#7);
  \draw[guide] (#5,#2-0.5,#7) -- (#5,#2-0.5,#3-0.5);
  \draw[guide] (#5,#2+0.5,#7) -- (#5,#2+0.5,#3-0.5);
  \draw[guide] (#5,#6,#3-0.5) -- (#1-0.5,#6,#3-0.5);
  \draw[guide] (#5,#6,#3+0.5) -- (#1-0.5,#6,#3+0.5);
  \draw[guide] (#5,#6,#3-0.5) -- (#5,#2-0.5,#3-0.5);
  \draw[guide] (#5,#6,#3+0.5) -- (#5,#2-0.5,#3+0.5);
}
%% \NewDocumentCommand{\ProjectVoxel}{m m m m O{\ProjX} O{\ProjY} O{\ProjZ}}{%
%%   \draw[fill=#4!10!white,draw=black,thin] (#1-0.5,#2-0.5,#7) -- (#1+0.5,#2-0.5,#7) -- (#1+0.5,#2+0.5,#7) -- (#1-0.5,#2+0.5,#7) -- cycle;
%%   \draw[fill=#4!30!white,draw=black,thin] (#5,#2-0.5,#3-0.5) -- (#5,#2-0.5,#3+0.5) -- (#5,#2+0.5,#3+0.5) -- (#5,#2+0.5,#3-0.5) -- cycle;
%%   \draw[fill=#4!40!white,draw=black,thin] (#1-0.5,#6,#3-0.5) -- (#1+0.5,#6,#3-0.5) -- (#1+0.5,#6,#3+0.5) -- (#1-0.5,#6,#3+0.5) -- cycle;
%% }

\begin{tikzpicture}[math3d,scale=0.8]
  %%[x={(1cm,-.35cm)},y={(.9cm,.25cm)},z={(0cm,1cm)},line cap=round]
\def\h{0.5}
\def\R{2.2}
\def\Tmin{-2.0349560*2}
\def\Tmax{1.9103089*2}
\def\Px{11}
\def\Py{7}
\def\Pz{7}

% X(t)=(t,0.42t+0.0375t^2,0.20t-0.0275t^2)
% p=(0,0,0), r=4, h=0.5
% Cube centers in world coordinates
\def\CubeCenters{%
  -3.750/-1.250/-1.250,
  -3.250/-1.250/-1.250,
  -3.250/-0.750/-1.250,
  -3.250/-0.750/-0.750,
  -2.750/-0.750/-0.750,
  -2.250/-0.750/-0.750,
  -1.750/-0.750/-0.750,
  -1.750/-0.750/-0.250,
  -1.250/-0.750/-0.250,
  -1.250/-0.250/-0.250,
  -0.750/-0.250/-0.250,
  -0.250/-0.250/-0.250,
  -0.250/-0.250/0.250,
  -0.250/0.250/-0.250,
  -0.250/0.250/0.250,
  0.250/-0.250/-0.250,
  0.250/-0.250/0.250,
  0.250/0.250/-0.250,
  0.250/0.250/0.250,
  0.750/0.250/0.250,
  1.250/0.250/0.250,
  1.250/0.750/0.250,
  1.750/0.750/0.250,
  2.250/0.750/0.250,
  2.250/1.250/0.250,
  2.750/1.250/0.250,
  2.750/1.750/0.250,
  3.250/1.750/0.250,
  3.750/1.750/0.250
}

% Grid indices of the same cubes
\def\CubeIndices{%
%  -8/-3/-3,
  -7/-3/-3,
  -7/-2/-3,
  -7/-2/-2,
  -6/-2/-2,
  -5/-2/-2,
  -4/-2/-2,
  -4/-2/-1,
  -3/-2/-1,
  -3/-1/-1,
  -2/-1/-1,
  -1/-1/-1,
  -1/-1/0,
  -1/0/-1,
  -1/0/0,
  0/-1/-1,
  0/-1/0,
  0/0/-1,
  0/0/0,
  1/0/0,
  2/0/0,
  2/1/0,
  3/1/0,
  4/1/0,
  4/2/0,
  5/2/0,
  5/3/0,
  6/3/0%,  7/3/0
}

\begin{scope}
% Coordinate axes
\draw[->,gray!70] (-0.5,0,0)--(18.5,0,0) node[below] {$x_1$};
\draw[->,gray!70] (0,-0.5,0)--(0,10.5,0) node[right] {$x_2$};
\draw[->,gray!70] (0,0,-0.5)--(0,0,7.5) node[above] {$x_3$};
\ProjectVoxel{\Px-8}{\Py-3}{\Pz-3}{25}
\foreach \x/\y/\z in \CubeIndices {
  \ProjectVoxel{\Px+\x}{\Py+\y}{\Pz+\z}{50}
}
\ProjectVoxel{\Px+7}{\Py+3}{\Pz+0}{25}
\DrawVoxel{\Px-8}{\Py-3}{\Pz-3}{white!70!black}
\foreach \x/\y/\z in \CubeIndices {
  \DrawVoxel{\Px+\x}{\Py+\y}{\Pz+\z}{white!25!black}
}
\DrawVoxel{\Px+7}{\Py+3}{\Pz+0}{white!70!black}
\end{scope}
\def\Radius{8}

% Centre de la sphère dans l'espace
\coordinate (P) at (\Px,\Py,\Pz);


% Curve, restricted to the Euclidean ball B(p,R), p=(0,0,0).
% Parametric curve:
\begin{scope}
\draw[blue,very thick,domain=-4:4,samples=150,variable=\t]
  plot ({\Px+(\t)/\h},{\Py+(0.42*\t+0.0375*\t*\t)/\h},
        {\Pz+(0.20*\t-0.0275*\t*\t)/\h});
\draw[black,thick,domain=-6:-4,samples=50,variable=\t]
  plot ({\Px+(\t)/\h},{\Py+(0.42*\t+0.0375*\t*\t)/\h},
        {\Pz+(0.20*\t-0.0275*\t*\t)/\h});
%\draw[blue!75!black,very thick,samples=180,domain=\Tmin:\Tmax,smooth,variable=\t]
% plot ({\Px+\t},{\Py+0.42*\t+0.075*\t*\t},{\Pz+0.20*\t-0.055*\t*\t});
\fill (\Px,\Py,\Pz) circle (1.6pt) node[blue,above left] {$p$};
\node[blue,above left] at (\Px+7,\Py+3,\Pz+0) {$X\cap B(p,r)$};
\end{scope}
% Optional sphere guide (schematic outline)
% \draw[blue!40,dashed] (0,0,0) circle (\R cm);
% Dessin du contour apparent dans le plan de la feuille
\begin{scope}[x=0.4cm,y=0.4cm]
  \pgftransformreset
  \path (P);
  \pgfgetlastxy{\ScreenX}{\ScreenY}

  \shade[ball color=cyan!35, opacity=0.15]
    (\ScreenX,\ScreenY) circle[radius=\Radius];

  \draw[blue!60, thin]
    (\ScreenX,\ScreenY) circle[radius=\Radius];
\end{scope}
\draw[black,thick,domain=4:6,samples=50,variable=\t]
  plot ({\Px+(\t)/\h},{\Py+(0.42*\t+0.0375*\t*\t)/\h},
        {\Pz+(0.20*\t-0.0275*\t*\t)/\h});

\end{tikzpicture}

  \end{center}
  \caption{Local dimension analysis on a voxelization of a 3D curve
    $X$ using the stability of measures $W_m$. Here $Y:=U_1(X) \cap
    B(p,r)$ with $r=8$, Measuring 1D projections gives
    $\HH^1(\pi_{1}(Y))=14+2\frac{1}{2}$,
    $\HH^1(\pi_{2}(Y))=5+2\frac{1}{2}$, $\HH^1(\pi_{3}(Y))=4$ (to
    simplify, we have counted $\frac{1}{2}$ for cubes on the boundary
    of the ball $\partial B(p,r)$ in order to approximate the
    intersections with the ball). So $\frac{1}{r^2}W_1(Y) =
    \frac{15^2+6^2+4^2}{8^2} \approx 4.32$. For a 1-manifold, expected
    value is $\bar{W}_1(1)=4$ while for a 2-manifold it is
    $\bar{W}_1(2)=8$. So the first measure $W_1$ seems to indicate
    that $p$ is part of a 1-manifold. Now measuring 2D projections
    gives $\HH^2(\pi_{1,2}(Y))=21+2\frac{1}{2}$,
    $\HH^2(\pi_{1,3}(Y))=18+2\frac{1}{2}$,
    $\HH^2(\pi_{2,3}(Y))=10+\frac{1}{2}$. So $\frac{1}{r^4}W_2(Y) =
    \frac{22^2+19^2+10.5^2}{8^4} \approx 0.23$. For a 1-manifold,
    expected value is $\bar{W}_2(1)=0$, and for a 2-manifold it is
    $\bar{W}_2(2)=\pi^2$. The second measure $W_2$ tells that $p$ is
    not part of 2-manifold, which confirms that $p$ is indeed part of
    a curve in 3D.\label{fig-illustration}}
\end{figure}

Note that both measures $V_J$ and $W_m$ are also stable with respect
to small perturbations in the position of $p$. The result above shows
that tangent perturbations do not change the value $W_m$. Although we
do not prove it here due to space reasons, a perturbation of $\nu$ in
the normal direction induces an error in $W_m$ of $O(\nu^2
r^{2m-2}+\nu r^{2m} / \rho)$. The following subsections give the
proofs of the theorems mentioned above.

\subsection{Local measures of $k$-balls are invariants}

The fact that the Pythagorean theorem extends to arbitrary volumes has
been rediscovered under different forms several times. The recent
paper of Ancel \cite{ancel2023pythagorean} provides a general form
for a Lebesgue measurable set $S$ in a $k$-dimensional affine
space. Then $\HH^k(S)^2=\sum_{J \in \mathcal{I}_k} \HH^k(\pi_J(S))^2$ (here
Lebesgue and $k$-th Hausdorff measures are the same).

We provide here a much simpler proof for the case of the ball, which
is sufficient in our context. This formulation also extends the result to lower dimensional projections.
\begin{theorem}\label{thm:ball-invariants}
Let $B^k$ be the unit Euclidean ball embedded isometrically into
$\mathbb{R}^d$ with the linear map $A\in\mathbb{R}^{d\times k}, A^T
A=\mathrm{Id}_k$ (so $k \le d$).  For every integer $m$, $0\le m\le k$, one has
\begin{equation}
  \label{eq:ball-invariants}
  \sum_{J \in \mathcal{I}_m} \HH^m(\pi_J(B^k))^2 = \omega_m^2 \binom{k}{m}.
\end{equation}
The same sum is obviously null if $k < m \Le d$.
\end{theorem}
\begin{proof}
For every subset $J \in \mathcal{I}_m$, Lemma~\ref{lem:svd} (see appendix) applied
to $L=\pi_J\circ A$ gives
\begin{equation}
  \label{eq-vol-proj-ball-1}
  \HH^m(\pi_J( B^k ))^2
  =\omega_m^2 \det(A_J A_J^T),
\end{equation}
where $A_J$ is the matrix containing the rows of $A$ indexed by
$J$. We denote by $\mathcal{I}^k_m$ the set of subsets of
$\{1,\dots,k\}$ with size $m$ (remember that $A_J$ may not be a square
matrix).  By the Cauchy--Binet formula, we have
\begin{equation}
  \label{eq-cauchy-binet-1}
\det(A_J A_J^T)
=
\sum_{I \in \mathcal{I}^k_m}
\det(A_J^I) \det( (A_J^T)_I ),
=
\sum_{I \in \mathcal{I}^k_m}
\det(A_J^I)^2,
\end{equation}
where $A_J^I$ denotes the $m\times m$ submatrix obtained by selecting
the rows indexed by $J$ and the columns indexed by $I$, and where we
have used that $(A_J^T)_I = (A_J^I)^T$ and the equality of the
determinant of a matrix and its transpose.
Therefore, \Equ{eq-vol-proj-ball-1} and \Equ{eq-cauchy-binet-1} implies that
\begin{align}
  \sum_{J \in \mathcal{I}_m}
  \HH^m(\pi_J(B^k))^2
  & =
  \omega_m^2
  \sum_{J \in \mathcal{I}_m}
  \sum_{I \in \mathcal{I}^k_m}
  \det(A_J^I)^2
  \nonumber \\
  \label{eq-proj-vol-ball-2}
  & =
  \omega_m^2
  \sum_{I \in \mathcal{I}^k_m}
  \sum_{J \in \mathcal{I}_m}
  \det(A_J^I)^2.
  && \text{(swapping the sums)}
\end{align}
Fix a subset $I$. The matrix $A^I$ consisting of the corresponding $m$
columns of $A$ satisfies $(A^I)^T A^I=\mathrm{Id}_m$, because the
columns of $A$ are orthonormal.
Applying Cauchy--Binet once more, we get
\begin{equation}
  \label{eq-cauchy-binet-2}
  1 = \det(\mathrm{Id}_m) = \det((A^I)^T A^I) = \sum_{J \in \mathcal{I}_m} \det(A_J^I)^2.
\end{equation}
Since there are $\binom{k}{m}$ possible subsets $I$ in $\mathcal{I}^k_m$, inserting \Equ{eq-cauchy-binet-2} into \Equ{eq-proj-vol-ball-2} gives

\[
\sum_{J \in \mathcal{I}_m}
\HH^m(\pi_J(B^k))^2
=
\omega_m^2
\binom{k}{m}.
\]

This completes the proof. \qed
\end{proof}
Thus the result is independent of the orientation of the ball in the
embedding space, as well as the dimension of this space. Let us denote
$\bar{W}_{m}(k)$ these geometric invariants of the $k$-ball. So for $0
\Le m \Le k$, $\bar{W}_{m}(k)=\omega_m^2\binom{k}{m}$ and for $k < m
\Le d$, $\bar{W}_{m}(k)=0$.  An immediate corollary is that, for $T$
any affine $k$-dimensional plane and $p \in T$, $W_m( B_T(p,r),p,r) =
\bar{W}_{m}(k) r^{2m}$.  These values (scaled by $r^{2m}$) correspond
to the expected values of the local measures when analyzing
digitizations of $k$-dimensional submanifolds (see appendix for tables
expliciting their values in low dimension).





\subsection{Local measures are stable in a neighborhood of a submanifold}

For a positive value $\delta$, we use the notation $(Y)_\delta$ to
denote the dilation of $Y$ with the ball of radius $\delta$. The
tangent space $T_pX$ of $X$ at $p$ is simply written $T$, and we
identify the affine tangent plane with $p+T$. We write simply
$D_r:=p+(T\cap B(0,r))$ for the disk of radius $r$ placed on the
tangent plane at $p$.

\begin{figure}[t]
  
\def\Angle{-28}
\def\Curvature{0.24}
\def\Radius{2.05}
\def\Mesh{0.30}
%\def\Eta{0.15} % Boundary-layer half-width is 2\eta
\def\Px{0}
\def\Py{1.65}
\def\Shift{0.75}
\def\Epsilon{1.5*\Mesh}
%\def\Eta{\Epsilon+\Radius*\Radius*\Curvature}
\def\Eta{0.2}

\pgfmathsetmacro{\ct}{cos(\Angle)}
\pgfmathsetmacro{\st}{sin(\Angle)}
\pgfmathsetmacro{\umin}{-2.9}
\pgfmathsetmacro{\umax}{2.9}
\pgfmathsetmacro{\uX}{0.9*\umax}
\pgfmathsetmacro{\rr}{\Radius-0.18}
% Endpoint of the curve inside the ball, computed numerically by bisection.
\pgfmathsetmacro{\ub}{\Radius}
\foreach \i in {1,...,25}{\pgfmathsetmacro{\ub}{0.5*(\ub+\Radius/sqrt(1+0.25*\Curvature*\Curvature*\ub*\ub))}}
\pgfmathsetmacro{\leftX}{\Px-\ub*\ct-0.5*\Curvature*\ub*\ub*\st}
\pgfmathsetmacro{\rightX}{\Px+\ub*\ct-0.5*\Curvature*\ub*\ub*\st}
\pgfmathsetmacro{\ellX}{\Px-\Radius*\ct}
\pgfmathsetmacro{\rX}{\Px+\Radius*\ct}
\pgfmathsetmacro{\rminus}{\rr*\ct}
%% \pgfmathsetmacro{\fleft}{\leftX-0.11}
%% \pgfmathsetmacro{\fright}{\rightX+0.11}
\pgfmathsetmacro{\fleft}{\leftX-0.6*\Mesh}
\pgfmathsetmacro{\fright}{\rightX-0.1*\Mesh}
\pgfmathsetmacro{\band}{2*\Eta}
\tikzset{curve/.style={black,very thick}, tangent/.style={red!80!black,dashed,thick}, ball/.style={blue!75!black,dashed,thick}, guide/.style={gray!65,densely dotted}, axis/.style={-Latex,thin}}
\begin{tikzpicture}[scale=1]
%%\foreach \shift/\panel in {0/{(a) Local projections},8.2/{(b) Homotopy}}{
\foreach \shift in {0,6.2}{
\begin{scope}[xshift=\shift cm]
  \draw[axis] (-2.6,0)--(2.75,0) node[right] {$x$};
  \draw[axis] (-2.5,-0.2)--(-2.5,3.5) node[above] {$y$};
  %\node[anchor=west,font=\bfseries] at (-3,4.75) {\panel};
  \draw[ball] (\Px,\Py) circle[radius=\Radius];
  \node[blue!70!black] at (1.65,3.55) {$B(p,r)$};
  % Voisinage epsilon de X ∩ B(p,r), tronqué par B(p,r)
  \draw[tangent] ({\Px+\umin*\ct},{\Py+\umin*\st})--({\Px+\umax*\ct},{\Py+\umax*\st});
  \draw[curve,domain=\umin:\umax,samples=100,smooth,variable=\u] plot ({\Px+\u*\ct-0.5*\Curvature*\u*\u*\st},{\Py+\u*\st+0.5*\Curvature*\u*\u*\ct});
  \node[anchor=south] at ({\Px+\uX*\ct-0.5*\Curvature*\uX*\uX*\st},{\Py+\uX*\st+0.5*\Curvature*\uX*\uX*\ct}) {$X$};
  \fill (\Px,\Py) circle (1.8pt) node[above right] {$p$};
  \node[red!75!black] at (2.45,0.72) {$T_pX$};
\end{scope}
}
% Voxel squares: a grid-based approximation, clipped visually by the ball.
\begin{scope}
\clip (\Px,\Py) circle[radius=\Radius];
\foreach \i in {-10,...,10}{
 \foreach \j in {-4,...,16}{
  \pgfmathsetmacro{\xx}{(\i+0.5)*\Mesh}
  \pgfmathsetmacro{\yy}{(\j+0.5)*\Mesh}
  \pgfmathsetmacro{\uu}{(\xx-\Px)*\ct+(\yy-\Py)*\st}
  \pgfmathsetmacro{\vv}{-(\xx-\Px)*\st+(\yy-\Py)*\ct}
  \pgfmathsetmacro{\distance}{abs(\vv-0.5*\Curvature*\uu*\uu)}
  \ifdim\distance pt<0.22pt
    \filldraw[fill=blue!25,draw=blue!40,line width=.18pt] (\i*\Mesh,\j*\Mesh) rectangle ++(\Mesh,\Mesh);
  \fi
 }
}
\end{scope}
% Redraw curve and tangent for clarity
\draw[tangent] ({\Px+\umin*\ct},{\Py+\umin*\st})--({\Px+\umax*\ct},{\Py+\umax*\st});
\draw[curve,domain=\umin:\umax,samples=100,variable=\u] plot ({\Px+\u*\ct-0.5*\Curvature*\u*\u*\st},{\Py+\u*\st+0.5*\Curvature*\u*\u*\ct});
\fill (\Px,\Py) circle (1.8pt);

% Plot dilated set
\begin{scope}
  \clip (\Px,\Py) circle (\Radius);
  \draw[
    magenta,
    opacity=0.18,
    line width=2*\Epsilon cm,
    line cap=round,
    line join=round
  ]
  plot[domain=\umin:\umax, samples=200, variable=\u]
    ({\Px+\u*cos(\Angle)-0.5*\Curvature*\u*\u*sin(\Angle)},
     {\Py+\u*sin(\Angle)+0.5*\Curvature*\u*\u*cos(\Angle)});
\end{scope}

% Projected intervals: E tangent disk, F curve projection, blue enlarged voxel envelope.
\draw[guide] (\ellX,\Py-\Radius*\st)--(\ellX,-1.1+\Shift);
\draw[guide] (\rX,\Py+\Radius*\st)--(\rX,-1.1+\Shift);
\draw[guide] (\fleft,-1.1+\Shift)--(\fleft,\Py-\Radius*\st+0.5*\Mesh);
\draw[guide] (\fright,-1.1+\Shift)--(\fright,\Py+\Radius*\st+1.7*\Mesh);
\draw[red!75!black,line width=3pt,opacity=.65] (\ellX,-.35+\Shift)--(\rX,-.35+\Shift);
\node[red!75!black,left,font=\scriptsize] at (\ellX,-.33+\Shift) {$\mathcal E_J$};
\draw[blue!75!black,line width=3pt,opacity=.65] (\fleft,-.6+\Shift)--(\fright,-.6+\Shift);
\node[blue!75!black,left,font=\scriptsize] at (\fleft,-0.6+\Shift) {$\mathcal F_J$};
% Boundary of E_J thickened by 2 eta: two neighborhoods of its endpoints.
\foreach \edge in {\ellX,\rX}{
  \draw[green!45!black,opacity=.48,line width=3pt,opacity=.65] ({\edge-\band},-1.1+\Shift) -- ({\edge+\band},-1.1+\Shift);
  %%\fill[green!45!black,opacity=.48] ({\edge-\band},-1.25) rectangle ({\edge+\band},-1.15);
  %%\draw[green!45!black,thin] ({\edge-\band},-1.75) rectangle ({\edge+\band},-1.52);
}
\node[green!35!black,font=\scriptsize,fill=white] at (0.5*\ellX+0.5*\rX,-1.1+\Shift)
 {$(\partial\mathcal E_J)_{2\eta}$};
% Symmetric difference: pieces of E_J not in F_J and F_J not in E_J.
% These are calculated from the same endpoints used for the intervals above.
\pgfmathsetmacro{\missingLeft}{max(0,\fleft-\ellX)}
\pgfmathsetmacro{\extraLeft}{max(0,\ellX-\fleft)}
\pgfmathsetmacro{\missingRight}{max(0,\rX-\fright)}
\pgfmathsetmacro{\extraRight}{max(0,\fright-\rX)}
\ifdim\missingLeft pt>0pt
 \draw[orange!85,line width=3pt,opacity=.65] (\ellX,-0.85+\Shift) -- (\fleft,-0.85+\Shift);
% \fill[orange!85] (\ellX,-1.0) rectangle (\fleft,-0.89);
\fi
\ifdim\extraLeft pt>0pt
 \draw[orange!85,line width=3pt,opacity=.65] (\fleft,-0.85+\Shift) -- (\ellX,-0.85+\Shift);
% \fill[orange!85] (\fleft,-1.0) rectangle (\ellX,-0.89);
\fi
\ifdim\missingRight pt>0pt
 \draw[orange!85,line width=3pt,opacity=.65] (\fright,-0.85+\Shift) -- (\rX,-0.85+\Shift);
%% \fill[orange!85] (\fright,-1.0) rectangle (\rX,-0.89);
\fi
\ifdim\extraRight pt>0pt
 \draw[orange!85,line width=3pt,opacity=.65] (\rX,-0.85+\Shift) -- (\fright,-0.85+\Shift);
% \fill[orange!85] (\rX,-1.0) rectangle (\fright,-0.89);
\fi
\node[orange!75!black,font=\scriptsize] at (0.5*\ellX+0.5*\rX,-0.87+\Shift)
 {$\mathcal E_J\triangle\mathcal F_J$};
%% \node[anchor=west,font=\scriptsize] at (-2.9,-2.15)
%%  {$\mathcal E_J\triangle\mathcal F_J\subset(\partial\mathcal E_J)_{2\eta}$};
% Second panel: homotopy curves use SAME angle and curvature, all tangent at p.
\begin{scope}[xshift=6.2cm]
\foreach \t in {0.25,0.5,0.75}{
\draw[gray!70,line width=.9pt,domain=-\rr:\rr,samples=70,variable=\u] plot ({\Px+\u*\ct-\t*0.5*\Curvature*\u*\u*\st},{\Py+\u*\st+\t*0.5*\Curvature*\u*\u*\ct});
}
\draw[very thick,orange!85!black,domain=-\rr:\rr,samples=70,variable=\u] plot ({\Px+\u*\ct-0.5*\Curvature*\u*\u*\st},{\Py+\u*\st+0.5*\Curvature*\u*\u*\ct});
\draw[very thick,red!70!black] ({\Px-\rr*\ct},{\Py-\rr*\st})--({\Px+\rr*\ct},{\Py+\rr*\st});
\node[anchor=west,fill=white,inner sep=1pt,font=\small] at (-2.3,.25) {$G_{J,t}(u)=p_x+u\cos\theta-tf(u)\sin\theta$};
\node[gray!70!black] at (2.0,3.1) {$0<t<1$};
\node[orange!80!black] at (2.1,2.8) {$t=1$};
\node[red!75!black] at (2.1,2.5) {$t=0$};

% Intervals
% Projection direction
\draw[-Latex,gray!70] (\Px,0.8*\Py)--(\Px,0.5*\Py) node[midway,right,font=\scriptsize] {$\pi_J$};
%\node[anchor=west,font=\scriptsize] at (-2.95,-0.40) {$G_{J,t}(u)=p_x+u\cos\theta-tf(u)\sin\theta$};
% Each interval has endpoints G_{J,t}(-r_-) and G_{J,t}(r_-).
% Vertical offsets are diagrammatic, not part of the codomain.
\foreach \t/\yy/\col in {0/0/red!75!black,0.25/-0.1/gray!65,0.5/-0.2/gray!65,0.75/-0.3/gray!65,1/-0.4/orange!85!black}{
  \pgfmathsetmacro{\aa}{\Px-\rr*\ct-\t*0.5*\Curvature*\rr*\rr*\st}
  \pgfmathsetmacro{\bb}{\Px+\rr*\ct-\t*0.5*\Curvature*\rr*\rr*\st}
  \draw[\col,line width=2.6pt,line cap=round] (\aa,\yy)--(\bb,\yy);
  \fill[\col] (\aa,\yy) circle (1.5pt) (\bb,\yy) circle (1.5pt);
}
\node[anchor=east,font=\scriptsize,red!75!black] at (-1.6,0) {$t=0$};
\node[anchor=east,font=\scriptsize,gray!65] at (-1.6,-0.2) {$0<t<1$};
\node[anchor=east,font=\scriptsize,orange!85!black] at (-1.6,-0.4) {$t=1$};
%% \node[anchor=north,font=\scriptsize] at (0,-2.30)
%%  {$G_{J,t}(B_{L_J}(0,r_-))\subset E_J\simeq\mathbf R$};
%% \node[anchor=north,font=\tiny,gray!70!black] at (0,-2.62)
%%  {Vertical offsets indicate $t$ only; all images lie on the $x$-axis.};

\end{scope}
\end{tikzpicture}

  \caption{\label{fig-proof-stability}Illustration of the proof of \RefTheorem{thm:stabilite-projections}. \textbf{Left}: the subset $Y$ is any covering set of $X \cap B(p,r)$ and covered by $X_\varepsilon \cap B(p,r)$, examplified here as the $h$-voxelization of $X$; Projections $\mathcal{E}_J$ and $\mathcal{F}_J$ are close. \textbf{Right}: the homotopy $G_{J,t}$ between subsets of $T \cap B(p,r)$ and $X \cap(B,r)$ (vertical offsets indicate $t$ only, all images lie on the $x$-axis).  }
\end{figure}

\begin{theorem}
\label{thm:stabilite-projections}
Under the hypotheses of \RefTheorem{thm:main}, there exist constants
$\alpha_0>0$ and $C>0$, both dependent solely on $d$ and $k$, such that, for
$0<r\Le \alpha_0\rho$ and $0\Le\varepsilon\Le r$, for every compact
set $Y\subset\mathbb R^d$ satisfying
\begin{equation}
\label{eq:encadrement-Y}
    X\cap B(p,r) \subset Y
    \subset (X\cap B(p,r))_\varepsilon\cap B(p,r),
\end{equation}
then, for every $J\in\mathcal I_m$,
\begin{align}
\label{eq:stabilite-VJ}
\left|V_J(Y,p,r)-V_J(D_r,p,r)\right|
&\Le C r^{m-1}\left(\varepsilon+\frac{r^2}{\rho}\right),\\
\label{eq:stabilite-W}
\left|W_m(Y,p,r)-\omega_m^2\binom{k}{m}r^{2m}\right|
&\Le C r^{2m-1}\left(\varepsilon+\frac{r^2}{\rho}\right).
\end{align}
\end{theorem}
\begin{proof}
To simplify expressions, $C$ denotes a positive constant depending only
on $d$ and $k$, whose value may change during the proof. Set
$X_r:=X\cap B(p,r)$.
\begin{enumerate}
  %% 1 %%
\item The proof is illustrated on \RefFigure{fig-proof-stability}. We start by approximating $X_r$ by $D_r$. Since $X$ has positive
reach, for a sufficiently small $\alpha_0$, the set $X$ is, in a
neighborhood of $p$, the graph over $T$ of a $C^2$ map
$f:\Omega\subset T\to T^\perp$, with
\(f(0)=0, Df(0)=0, |f(u)|\Le C\frac{|u|^2}{\rho}, \|Df(u)\|\Le C\frac{|u|}{\rho}.\)
Let $\Phi(u):=p+u+f(u)$. For $x=\Phi(u)\in X_r$, since
$u\perp f(u)$, we have $|u|\Le |x-p|\Le r$. Furthermore,
\(|x-(p+u)|=|f(u)|\Le C\frac{r^2}{\rho}=: \delta,\)
so that
\begin{equation}
\label{eq:X-outer}
 X_r\subset (D_r)_\delta.
\end{equation}
Reciprocally, one can choose
\(r_-:=r-C\frac{r^3}{\rho^2}>0,\)
with $\alpha_0$ small enough so that $r-r_-\Le\delta$, and such that
\begin{equation}
\label{eq:inner-graph}
 \Phi(B_T(0,r_-))\subset X_r.
\end{equation}
Indeed, for $|u|\Le r_-$,
\(|\Phi(u)-p|^2=|u|^2+|f(u)|^2 \Le |u|^2+C\frac{|u|^4}{\rho^2},\)
and the choice of $r_-$ ensures that the right-hand side is at most
$r^2$.

\item We now compare the projection of $Y$,
$\mathcal F_J:=\pi_J(Y)$, with the projected tangent disk
$\mathcal E_J:=\pi_J(D_r)$. Notice that $Y\subset B(p,r)$ by
\eqref{eq:encadrement-Y}; hence
\(V_J(Y,p,r)=\HH^m(\mathcal F_J).\)
Using \eqref{eq:encadrement-Y} and the fact that $\pi_J$ is
$1$-Lipschitz, we obtain
\(\pi_J(X_r)\subset\mathcal F_J \subset \pi_J(X_r)+B^m_\varepsilon.\)
Together with \eqref{eq:X-outer}, this yields
\begin{equation}
\label{eq:outer-projection}
 \mathcal F_J\subset\mathcal E_J+B^m_\eta,
 \qquad
 \eta:=\varepsilon+\delta
 =\varepsilon+C\frac{r^2}{\rho}.
\end{equation}
Since $\varepsilon\Le r$ and $r\Le\alpha_0\rho$, we have
$\eta\Le C r$.

\item Consider the linear map
\(A_J:=\pi_J|_T:T\to\R^m.\)
Let $1\Ge\sigma_1\Ge\cdots\Ge\sigma_m\Ge0$ be its $m$ largest
singular values. The semi-axes of $\mathcal E_J$ are
$r\sigma_1,\ldots,r\sigma_m$, and therefore
\begin{equation}
\label{eq:ellipsoid-volume}
 \HH^m(\mathcal E_J)
 =\omega_m r^m\prod_{i=1}^m\sigma_i.
\end{equation}
Fix from now on a sufficiently large numerical constant $C_0>2$,
independent of $J,r,\rho$ and $\varepsilon$, and distinguish two cases.

\item Assume first that
\(r\sigma_m\Le C_0\eta.\)
Then $A_J$ is almost degenerate. In principal coordinates,
\(\mathcal E_J\subset\prod_{i=1}^m[-r\sigma_i,r\sigma_i].\)
The first $m-1$ dimensions are at most $2r$, whereas the last one is
at most $2C_0\eta$. Hence the $m$-volume of $\mathcal E_J$ is upper bounded as
%\begin{equation}
%\label{eq:degenerate-ellipsoid-volume}
 $\HH^m(\mathcal E_J)\Le C r^{m-1}\eta$.
%\end{equation}

Moreover, by \eqref{eq:outer-projection},
$\mathcal F_J\subset\mathcal E_J+B^m_\eta$, which is contained in
\(\prod_{i=1}^m[-r\sigma_i-\eta,r\sigma_i+\eta].\)
Since $\eta\Le Cr$, the first $m-1$ dimensions have size $O(r)$ and
the last one has size $O(\eta)$. Thus the $m$-volume of $\mathcal F_J$ is upper bounded similarly as
%\begin{equation}
%\label{eq:degenerate-F-volume}
$\HH^m(\mathcal F_J)\Le C r^{m-1}\eta$.
%\end{equation}

Combining the last two estimates gives
\begin{equation}
\label{eq:degenerate-result}
 \left|\HH^m(\mathcal F_J)-\HH^m(\mathcal E_J)\right|
 \Le C r^{m-1}\eta.
\end{equation}

\item Assume now that
\(r\sigma_m>C_0\eta.\)
The condition $C_0>2$ ensures in particular that the ellipsoid has a
nonempty core after removing a boundary layer of width $2\eta$.
Let
\(L_J:=(\ker A_J)^\perp\cap T.\)
Then $\dim L_J=m$ and
$\widetilde A_J:=A_J|_{L_J}$ is an isomorphism from $L_J$ onto
$\R^m$. For every $s>0$,
\(A_J(B_T(0,s))=\widetilde A_J(B_{L_J}(0,s)).\)
Consider
\(G_J:B_{L_J}(0,r_-)\to\R^m, G_J(u):=\pi_J(\Phi(u)) =\pi_J(p)+\widetilde A_Ju+\pi_J(f(u)).\)
By \eqref{eq:inner-graph},
\begin{equation}
\label{eq:G-image}
 G_J(B_{L_J}(0,r_-))\subset\pi_J(X_r)\subset\mathcal F_J.
\end{equation}
Introduce the homotopy
\(G_{J,t}(u):=\pi_J(p)+\widetilde A_Ju+t\,\pi_J(f(u)), t\in[0,1],\)
and the contracted ellipsoid
\(\mathcal E_J^-:=\pi_J(p)+\widetilde A_J(B_{L_J}(0,r_-)).\)
For every $u\in\partial B_{L_J}(0,r_-)$ and $t\in[0,1]$,
\(|G_{J,t}(u)-G_{J,0}(u)|\Le |f(u)|\Le\delta,\)
hence
\begin{equation}
\label{eq:homotopy-boundary}
 G_{J,t}(\partial B_{L_J}(0,r_-))
 \subset(\partial\mathcal E_J^-)_\delta.
\end{equation}
To show that the contracted ellipsoid is included in the projection of $Y$, it suffices to prove that it is continuously transported through the homotopy without touching the border. 
Let $y\in\mathcal E_J^-$ satisfy
$\operatorname{dist}(y,\partial\mathcal E_J^-)>\delta$. Then $y$
does not belong to
$G_{J,t}(\partial B_{L_J}(0,r_-))$ for any $t\in[0,1]$. Therefore the Brouwer degree
\(\deg(G_{J,t},B_{L_J}(0,r_-),y)\)
is well defined and independent of $t$. For $t=0$, $G_{J,0}$ is an
affine isomorphism, hence
\(\deg(G_{J,0},B_{L_J}(0,r_-),y)=\pm1.\)
By homotopy invariance,
\(\deg(G_J,B_{L_J}(0,r_-),y)=\pm1\ne0,\)
and therefore $y\in G_J(B_{L_J}(0,r_-))\subset\mathcal F_J$. We have
proved
\begin{equation}
\label{eq:inner-contracted}
 \mathcal E_J^-\setminus(\partial\mathcal E_J^-)_\delta
 \subset\mathcal F_J.
\end{equation}
The ellipsoids $\mathcal E_J^-$ and $\mathcal E_J$ are homothetic and,
since $\pi_J$ is $1$-Lipschitz,
\(d_H(\mathcal E_J^-,\mathcal E_J) \Le r-r_-\Le\delta.\)
By convexity, any point of $\mathcal E_J$ at distance greater than
$\delta+(r-r_-)$ from its boundary belongs to $\mathcal E_J^-$ and is
at distance greater than $\delta$ from $\partial\mathcal E_J^-$. Since
\(\delta+(r-r_-)\Le2\delta\Le2\eta,\)
\eqref{eq:inner-contracted} yields
\begin{equation}
\label{eq:inner-final}
 \mathcal E_J\setminus(\partial\mathcal E_J)_{2\eta}
 \subset\mathcal F_J.
\end{equation}

\item Combining the outer and inner inclusions gives
\(\mathcal E_J\setminus\mathcal F_J \subset(\partial\mathcal E_J)_{2\eta},\)
while, since $\mathcal E_J$ is closed and convex,
\(\mathcal F_J\setminus\mathcal E_J \subset(\mathcal E_J+B^m_\eta)\setminus\mathcal E_J \subset(\partial\mathcal E_J)_\eta.\)
Consequently,
\begin{equation}
\label{eq:symmetric-difference}
 \mathcal F_J\triangle\mathcal E_J
 \subset(\partial\mathcal E_J)_{2\eta}.
\end{equation}
We now estimate the volume of this boundary layer. If
$K\subset\R^m$ is a convex body contained in a ball of radius $r$, then
for every fixed $a>0$ there exists $C_{m,a}>0$ such that, for
$0<t\Le ar$,
\begin{equation}
\label{eq:convex-boundary-layer}
 \HH^m((\partial K)_t)\Le C_{m,a}\,t r^{m-1}.
\end{equation}
Indeed, the outer layer is controlled by the Steiner formula. For the
inner layer, define
\(K_{-s}:=\{x\in K:\operatorname{dist}(x,\partial K)\Ge s\}.\)
Then
\(\HH^m(K)-\HH^m(K_{-t}) =\int_0^t \mathrm{Area}(K_{-s})\,ds \Le t \mathrm{Area}(K) \Le m\omega_m t r^{m-1},\)
where we used monotonicity of the area for convex bodies. Since
$\eta\Le Cr$, the estimate \eqref{eq:convex-boundary-layer} applies to
$t=2\eta$, and therefore
\(\HH^m((\partial\mathcal E_J)_{2\eta}) \Le C\eta r^{m-1}.\)
Together with \eqref{eq:symmetric-difference}, this gives
\begin{equation}
\label{eq:nondegenerate-result}
 \left|\HH^m(\mathcal F_J)-\HH^m(\mathcal E_J)\right|
 \Le \HH^m(\mathcal F_J\triangle\mathcal E_J) \Le C\eta r^{m-1}.
\end{equation}

\item Combining \eqref{eq:degenerate-result} and
\eqref{eq:nondegenerate-result} gives, in all cases,
\(\left|\HH^m(\mathcal F_J)-\HH^m(\mathcal E_J)\right| \Le C\eta r^{m-1}.\)
Since $Y\subset B(p,r)$ and $D_r\subset B(p,r)$, this is exactly
\eqref{eq:stabilite-VJ}.

\item Finally, set
\(a_J:=\HH^m(\pi_J(Y)),\qquad b_J:=\HH^m(\pi_J(D_r)).\)
Since both projected sets are contained in an $m$-ball of radius $r$,
$a_J+b_J\Le Cr^m$. Hence
\(|a_J^2-b_J^2| =|a_J-b_J|(a_J+b_J) \Le C r^{2m-1}\left(\varepsilon+\frac{r^2}{\rho}\right).\)
Summing over $J\in\mathcal I_m$ gives
\[
\left|W_m(Y,p,r)-W_m(D_r,p,r)\right| \Le C r^{2m-1}\left(\varepsilon+\frac{r^2}{\rho}\right).\]
\end{enumerate}
By \RefTheorem{thm:ball-invariants},
\(W_m(D_r,p,r)=\omega_m^2\binom{k}{m}r^{2m},\)
which proves \eqref{eq:stabilite-W}.
\qed
\end{proof}


\subsection{Discretization error}

We now show that the $h$-voxelization of $X$, restricted to $B(p,r)$,
satisfies the hypotheses of the previous stability theorem.  The only
geometric ingredient is the following elementary estimate on thin
annuli of $X$. In this section, we can relax slightly the $C^2$-smoothness
constraint on $X$ and ask $X\subset\mathbb R^d$ to be a closed
$C^{1,1}$-submanifold without boundary, with
$\operatorname{reach}(X)\Ge\rho>0$.


\begin{lemma}[Radial annulus estimate]
  \label{lem:radial-annulus}
For any $p,x\in X$ and $0<r<|x-p|=s<\rho$, one has
\begin{equation}\label{eq:radial-annulus}
\operatorname{dist}\bigl(x,X\cap B(p,r)\bigr)
\Le
\frac{s-r}{\sqrt{1-s^2/(4\rho^2)}}
\Le \frac{2}{\sqrt3}(s-r).
\end{equation}
\end{lemma}
\begin{proof}
Let $g:X\setminus\{p\}\to\mathbb R$ be defined by
$g(z)=|z-p|$. The angle estimate of Federer
(see also \cite[Lemma~9]{lieutier2024manifolds})
gives, for $|z-p|<\rho$,
$\|\operatorname{proj}_{(T_zX)^\perp}(z-p)\|
\Le |z-p|^2/(2\rho)$.
Consequently,
$\|\nabla_X g(z)\|
=\|\operatorname{proj}_{T_zX}(z-p)\|/|z-p|
\Ge\sqrt{1-|z-p|^2/(4\rho^2)}$.


Starting at $x$, follow the unit-speed descending
gradient field $-\nabla_Xg/\|\nabla_Xg\|$ until the
first point $x_*$ with $g(x_*)=r$.
Along this unit-speed curve $\gamma:[0,t_*]\to X$,
the distance to $p$ decreases and
$-\frac{d}{dt}g(\gamma(t))
=\|\nabla_Xg(\gamma(t))\|
\Ge c_s:=\sqrt{1-s^2/(4\rho^2)}$.
Integrating from $0$ to $t_*$ yields
$s-r\Ge c_s t_*$.
Since $\gamma$ has unit speed,
$\operatorname{Length}(\gamma)=t_*
\Le(s-r)/c_s$.
As $x_*=\gamma(t_*)\in X\cap B(p,r)$,
the conclusion follows.
%% Along this curve $\gamma$, the distance to $p$
%% decreases and
%% $-\frac{d}{dt}g(\gamma(t))
%% =\|\nabla_Xg(\gamma(t))\|
%% \Ge\sqrt{1-s^2/(4\rho^2)}$.
%% Thus the length of $\gamma$ is at most
%% $(s-r)/\sqrt{1-s^2/(4\rho^2)}$.
%% Since $x_*\in X\cap B(p,r)$,
%% $\operatorname{dist}(x,X\cap B(p,r))
%% \Le |x-x_*|\Le\operatorname{Length}(\gamma)$,
%% which proves the claim.
\qed
\end{proof}


%% \begin{lemma}[Radial annulus estimate]
%% \label{lem:radial-annulus}
%% There exist constants $\alpha_0=\alpha_0(d,k)>0$ and $C>0$ such that,
%% if $0<r\Le\alpha_0\rho$, $0\Le s\Le\sqrt d\,r$, and
%% $x\in X$ satisfies
%% \[
%%  r\Le |x-p|\Le r+s,
%% \]
%% then
%% \begin{equation}
%% \label{eq:radial-annulus}
%%  \operatorname{dist}\bigl(x,X\cap B(p,r)\bigr)
%%  \Le C\bigl(|x-p|-r\bigr)
%%  \Le Cs.
%% \end{equation}
%% \end{lemma}
%% \begin{proof}
%% Choose $\alpha_0$ sufficiently small so that the graph representation
%% of $X$ over $T_pX$ is valid in $B(p,(1+\sqrt d)r)$. Write
%% \(x=\Phi(u)=p+u+f(u),\)
%% with
%% \(\|Df(v)\|\Le C\frac{|v|}{\rho}.\)
%% Consider the radial graph curve
%% \(\gamma(t):=\Phi(tu),\qquad 0\Le t\Le1,\)
%% and set $g(t):=|\gamma(t)-p|$. Since $g(0)=0$ and $g(1)\Ge r$, there
%% exists $t_*\in[0,1]$ such that $g(t_*)=r$. Set
%% $x_*:=\gamma(t_*)\in X\cap B(p,r)$.

%% Now
%% \(\gamma'(t)=u+Df(tu)[u].\)
%% Since $|u|\Le|x-p|\Le(1+\sqrt d)r$, decreasing $\alpha_0$ if
%% necessary gives
%% \(|\gamma'(t)|\Le C_1|u|.\)
%% Moreover,
%% \(g'(t) =\frac{\langle\gamma(t)-p,\gamma'(t)\rangle} {|\gamma(t)-p|} =\frac{t|u|^2+\langle f(tu),Df(tu)[u]\rangle} {\sqrt{t^2|u|^2+|f(tu)|^2}}.\)
%% Using the standard graph estimates and again taking $\alpha_0$
%% sufficiently small, we obtain
%% \(g'(t)\Ge c|u|\)
%% for $t\in[t_*,1]$, with $c>0$ independent of $p,r$ and $x$.
%% Therefore
%% \(|x-p|-r =\int_{t_*}^1g'(t)\,dt \Ge c|u|(1-t_*),\)
%% whereas
%% \(|x-x_*| \Le\int_{t_*}^1|\gamma'(t)|\,dt \Le C_1|u|(1-t_*).\)
%% Hence
%% \(\operatorname{dist}(x,X\cap B(p,r)) \Le |x-x_*| \Le \frac{C_1}{c}\bigl(|x-p|-r\bigr),\)
%% which proves the claim.
%% \qed
%% \end{proof}

\begin{theorem}[Local cubical approximation]
  \label{thm:cubical-admissible-set}
%% Let $X\subset\mathbb R^d$ be a closed $C^{1,1}$
%% submanifold without boundary, with
%% $\operatorname{reach}(X)\Ge\rho>0$.
Let $U_h(X)$ be the union of all closed cubes of the
standard cubical grid of mesh size $h$ that intersect $X$.
For $p\in X$, set $X_r=X\cap B(p,r)$ and
$Y_h=U_h(X)\cap B(p,r)$.
If
\(
    0<h\Le r<\frac{\rho}{1+\sqrt d},
\)
then
\begin{equation}\label{eq:cubical-approximation}
    X_r\subset Y_h
    \subset (X_r)_{\varepsilon_h}\cap B(p,r),
    \qquad
    \varepsilon_h=
    \left(1+\frac{2}{\sqrt3}\right)\sqrt d\,h.
\end{equation}
\end{theorem}

\begin{proof}
The first inclusion follows from $X\subset U_h(X)$ (it is a covering
of $X$). For the second inclusion, for any $y\in Y_h$, choose a cube
containing $y$ and intersecting $X$, and let $x\in X$ belong to that
cube. Then $|x-y|\Le\sqrt d\,h$ and $|x-p|\Le r+\sqrt d\,h\Le(1+\sqrt
d)r<\rho$.  If $x\in X_r$, then $\operatorname{dist}(y,X_r)\Le\sqrt
d\,h$.  Otherwise, the radial annulus estimate
(Lemma~\ref{lem:radial-annulus}) gives $\operatorname{dist}(x,X_r)
\Le\frac{2}{\sqrt3}(|x-p|-r) \Le\frac{2}{\sqrt3}\sqrt d\,h$.
Consequently, $\operatorname{dist}(y,X_r) \Le
|y-x|+\operatorname{dist}(x,X_r) \Le(1+2/\sqrt3)\sqrt d\,h$, which
concludes. \qed
\end{proof}


%% \begin{theorem}[Local cubical approximation]
%% \label{thm:cubical-admissible-set}
%% Let $X_r:=X\cap B(p,r), Y_h:=U_h(X)\cap B(p,r)$.
%% There exist a constant $C=C(d,k)>0$
%% such that, if
%% %% There exist constants $\alpha_0=\alpha_0(d,k)>0$ and $C=C(d,k)>0$
%% %% such that, if
%% \(
%%  0< r \Le \frac{1}{1+\sqrt{d}}\rho, %% \alpha_0\rho,
%%  0<h\Le r,
%% \)
%% then
%% \begin{equation}
%% \label{eq:cubical-admissible}
%%  \boxed{
%%  X_r\subset Y_h
%%  \subset (X_r)_{\varepsilon_h}\cap B(p,r),
%%  \qquad
%%  \varepsilon_h=C\sqrt d\,h.}
%% \end{equation}
%% \end{theorem}
%% \begin{proof}
%% The first inclusion is immediate since the voxelization covers $X$.
%% Moreover, $Y_h\subset B(p,r)$ by definition.

%% Let $y\in Y_h$. Then $y$ belongs to a voxel $Q=z+Q_h$ such that
%% $Q\cap X\ne\emptyset$. Choose $x\in Q\cap X$. Since $Q$ has diameter
%% $\sqrt d\,h$,
%% \begin{equation}
%% \label{eq:y-x-cube}
%%  |y-x|\Le\sqrt d\,h.
%% \end{equation}
%% Also,
%% \(|x-p|\Le |x-y|+|y-p|\Le r+\sqrt d\,h.\)
%% If $x\in X_r$, then
%% $\operatorname{dist}(y,X_r)\Le\sqrt d\,h$. Otherwise,
%% \RefLemma{lem:radial-annulus}, applied with $s=\sqrt d\,h$, gives
%% \(\operatorname{dist}(x,X_r)\Le C\sqrt d\,h.\)
%% Thus in all cases
%% \(\operatorname{dist}(y,X_r) \Le |y-x|+\operatorname{dist}(x,X_r) \Le C'\sqrt d\,h.\)
%% Absorbing $C'$ into $C$ proves
%% \(Y_h\subset(X_r)_{\varepsilon_h}\cap B(p,r), \qquad \varepsilon_h=C\sqrt d\,h.\)
%% \qed
%% \end{proof}
